\documentclass[noauthor]{ximera}

\title{Week 6}
\author{Math 1193 Design Team}

\begin{document}
\begin{abstract}
    Solutions for Week 6.
\end{abstract}
\maketitle

\begin{problem}
Look over the following table classifying expressions.

\begin{center}
\begin{tabular}{|c|c|c|c|} 
 \hline
 Expression & Polynomial? & Degree & Leading Coefficient \\
 \hline\hline
 \(x\) & Yes & 1 & 1 \\ 
 \hline
 \(3x^2\) & Yes & 2 & 3 \\
 \hline
 \(7x + x^{1/2}\) & No & N/A & N/A \\
 \hline
 \(2x^3 - x^2 + x - 4\) & Yes & 3 & 2 \\
 \hline
 \(\pi x^2 - 9\) & Yes & 2 & \(\pi\) \\
 \hline
 \(\frac{2}{3}x^5 + \frac{9}{7}x\) & Yes & 5 & \(\frac{2}{3}\) \\
 \hline
 \(\frac{2x^5 + 9x}{3x^5 - 7x}\) & No & N/A & N/A \\
 \hline
 \(x^\pi - 2\) & No & N/A & N/A \\
 \hline
 \(2^x + 2^x + 1\) & No & N/A & N/A \\
 \hline
 \((2 - x)(x + 1)\) & Yes & 2 & \(-1\) \\
 \hline
\end{tabular}
\end{center}

\begin{enumerate}
    \item Come up with a definition of a polynomial. Reviewing the
    definition of a quadratic may help you. Check with your group
    and the instructor to confirm.
    \item Identify all of the following which are polynomials 
    and give their degree and leading coefficient.
    \begin{enumerate}
        \item \(f\left(x\right)=2x^3-x+1\)
	    \item \(g\left(x\right)=2x^\pi-x+1\)
	    \item \(h\left(x\right)=7-x\)
	    \item \(f\left(x\right)=x\)
	    \item \(g\left(x\right)=\left(x-17\right)\left(x+8\right)\)
	    \item \(h\left(x\right)=\left(2x-3\right)^5\left(x-1\right)^2\)
	    \item \(f\left(x\right)= e x^5-3\)
	    \item \(g\left(x\right)=\frac{x-1}{x+1}\)
	    \item \(h\left(x\right)=7x^6-2x+x^6\)
	    \item \(f\left(x\right)=1+x+x^2\)
	    \item \(g\left(x\right)=-3\left(x-\frac{1}{2}\right)^2+4\)
	    \item \(h\left(x\right)=\left(x^2+1\right)\left(x-1\right)\left(x+1\right)^2\)
	    \item \(f\left(x\right)=\left(-3x+2\right)\left(5x+1\right)\left(2x-2\right)^2\) 
    \end{enumerate}
\end{enumerate}

\begin{explanation}
\begin{enumerate}
    \item Here's the definition from the College Algebra textbook.
    
    ``Let \(n\) be a whole number and \(a_n\), \(a_{n - 1}\), \(a_{n - 2}\), \(\ldots\), \(a_1\), \(a_0\) be real numbers,
    where \(a_n \ne 0\). Then a function defined by 
    $$
    f(x) = a_nx^n + a_{n - 1}x^{n - 1} + a_{n - 2}x^{n-2} + \cdots + a_1 x + a_0
    $$
    is called a \textbf{polynomial function of degree \(n\)}.''

    This is a kind of technical-looking definition. Another way of thinking of polynomials is that polynomials are all the expressions you get from
    taking whole number powers of \(x\), multiplying by real numbers, then adding them. The highest exponent that appears is called the degree of the polynomial. 

    Notice that we'll say 0 and 1 are whole numbers, so \(x^1\), which we usually write as \(x\) is a polynomial. Same with \(x^0\), which we usually write as 1.

    Notice that when you multiply polynomials, you get another polynomial, since you can always rewrite in the above form.

    \item 
    \begin{enumerate}
        \item \(2x^3 - x + 1\) is a polynomial of degree 3. The leadning coefficient is the coefficient attached to the highest power of \(x\), which in this case is 2.
        \item \(2x^\pi - x + 1\) is not a polynomial. The exponents on \(x\) must all be whole numbers, and \(\pi\) is not a whole number.
        \item \(7 - x\) is a polynomial. Since the highest power of \(x\) that appears is \(x\), which we can rewrite as \(x^1\), its degree is 1 and its leading coefficient is \(-1\). Notice that 
        the order in which we write the terms does not matter, since addition's order does not matter.
        \item \(x\) is a polynomial of degree 1 with leading coefficient 1.
        \item \((x - 17)(x + 8)\) is a polynomial, since we can multiply the two factors to obtain something that is a sum of whole number powers of \(x\) with real 
        number coefficients (try it!). When you do this, you'll see that the degree is 2 and leading coefficient is 1.
        \item \((2x - 3)^5(x - 1)^2\) is a polynomial, since exponentiation is repeated multiplication, and multiplication of two polynomials results in a polynomial.
        
        However, the degree and leading coefficient are a bit more difficult to find. You could do all the multiplications, but that would take a while.
        
        Instead, consider the following. It's a long digression, but it may be useful.
        
        When we multiply something like \((x - 1)(x + 2)\), we do what's sometimes called FOILing. We multiply the first terms, \(x \cdot x = x^2\),
        then the outer terms, \(x \cdot 2 = 2x\), then the inner terms, \(-1 \cdot x = -x\), and finally the last terms, \(-1 \cdot 2 = -2\), then add them together,
        \(x^2 + 2x + (-x) + (-2) = x^2 + x - 2\).

        Although FOIL is a helpful way to remember this, this is a consequence of the distributive property.
        \begin{align*}
            (x - 1)(x + 2) & = (x - 1)x + (x - 1)2 \\
            & = x(x - 1) + 2(x - 1) \\
            & = x^2 - x + 2x - 2 \\
            & = x^2 + x - 2. 
        \end{align*}
        
        Another way to describe what's going on is that we're choosing one term from each set of parentheses and multiplying them, then doing the same for all possible
        choices, and finally adding our results, collecting like terms if need be.

        This works for other multiplication of polynomials as well! For example, \((x - 1)(x^2 + 2x + 1)\) cannot be done using FOIL (the \(2x\) in the second factor 
        isn't first, outer, inner, or last). Instead, we multiply the \(x\) in the first factor by \(x^2\), then \(2x\), then \(1\), then multiply the \(-1\) in the
        first factor by \(x^2\), then \(2x\), then \(1\). Notice how every possible choice of terms from each set of parentheses was multiplied together.

        Similarly, you can't do \((4x + 1)(x + 2)(x + 3)\) all at once with FOIL (what's the inner term of \((x + 2)\)?). You'd have to FOIL the first two factors, then 
        multiply by the last as shown above. 

        However, we can get information about the degree and leading coefficient without multiplying out the whole polynomial. Since the degree is the \textbf{highest}
        power of \(x\), the only choice of terms from each set of parentheses that gives you the highest power of \(x\) will be the highest power term in each set
        of parentheses. In the example of \((4x + 1)(x + 2)(x + 3)\), if we take the highest power term from each set of parentheses, we get \(4x \cdot x \cdot x = x^3\).
        No other choice gets you any power higher than \(x^2\). Therefore, \((4x + 1)(x + 2)(x + 3) = 4x^3 + \cdots\), where there aren't any other terms with an \(x^3\)
        in them. Thus, we know for sure that the degee of \((4x + 1)(x + 2)(x + 3)\) is 3 and its leading coefficient is 4.

        In the problem, we could rewrite \(2x - 3)^5(x - 1)^2\) as \((2x - 3)(2x - 3)(2x - 3)(2x - 3)(2x - 3)(x - 1)(x - 1)\) and choose the highest power 
        of \(x\) from each set of parentheses: \(2x \cdot 2x \cdot 2x \cdot 2x \cdot 2x \cdot x \cdot x = (2x)^5\cdot x^2 = 32x^7\). Since this is guaranteed
        to be the term with the highest power of \(x\), we know that the degree is 7 and the leading coefficient is 32.

        \item \(ex^5 - 3\) is a polynomial of degree 5 and leading coefficient \(e\). \(e\) is a very important real number constant you'll learn about later in this semester.
        
        \item \(\frac{x - 1}{x + 1}\) is not a polynomial. The quotient of two polynomials is not a polynomial.
        \item \(7x^6 - 2x + x^6\) is a polynomial. It can be rewritten as \(8x^6 - 2x\), so it is of degree 6 and has leading coefficient 8.
        \item \(1 + x + x^2\) is a polynomial of degree 2 and leading coefficient 1. Notice that the order of terms again does not matter when finding degree and leading coefficient.
        \item \(-3\left(x - \frac{1}{2}\right) + 4\) is a polynomial. This is the vertex form of a quadratic, which is a degree 2 polynomial. After we FOIL the 
        \(\left(x - \frac{1}{2}\right)^2\), we get an \(x^2\) term, then multiplying by \(-3\) gives us a leading coefficient of \(-3\).
        \item \((x^2 + 1)(x - 1)(x + 1)^2\) is a polynomial. Multiplying the leading terms of the factors gives us \(x^2 \cdot x \cdot x^2 = x^5\), so the degree is 5
        and the leading coefficient is 1.
        \item \((-3x + 2)(5x + 1)(2x - 2)^2\) is a polynomial. Multiplying the leading terms of the factors gives us \(-3x \cdot 5x \cdot (2x^2) = -60x^4\), so the degree is 4
        and the leading coefficient is \(-60\).
        

    \end{enumerate}
\end{enumerate}
\end{explanation}
\end{problem}

\begin{problem}
Solve the following polynomial equations.

\begin{enumerate}
    \item \((x - 1)(x + 2) = 0\)
    \item \((x - 1)^2(x + 2) = 0\)
    \item \(-3(x - 1)^2(x + 2) = 0\)
    \item \((2x - 1)^3(x + 1) = 0\)
    \item \(2x^3 + 8x^2 + 3x + 12 = 0\)
    \item \(2x^2 + 7x + 3 = 0\)
\end{enumerate}

How could you check your answers? 

Which equations were most difficult to solve?

\begin{explanation}
Remember that solving an equation means that we're finding which numbers, when plugged into the equation, make the equation true.

When numbers multiply to be 0, at least one of them has to be 0. Because of this, when setting expressions equal to 0 and solving,
we can set each factor equal to 0 and solve to obtain the solutions.
    \begin{enumerate}
        \item We set each factor equal to 0: \(x - 1 = 0\) and \(x + 2 = 0\), yielding \(x = -2, 1\).
        \item Again, we set each factor equal to 0: \(x - 1 = 0\), \(x - 1 = 0\), and \(x + 2 = 0\), yielding \(x = -2, 1\). The first two equations give the same solution.
        \item Same here: \(-3 = 0\), \(x - 1 = 0\), \(x - 1 = 0\), and \(x + 2 = 0\). The first equation is never true, so it doesn't give a solution. Therefore,
        the solutions are \(x = -2, 1\).
        \item Again: \(2x - 1 = 0\) and \(x + 1 = 0\), so \(x = \frac{1}{2}, -1\). Note that the power of 3 adds no solutions.
        \item The right-hand side of the equation isn't factored, so let's factor using grouping.
        \begin{align*}
            2x^3 + 8x^2 + 3x + 12 & = (2x^3 + 8x^2) + (3x + 12) \\
            & = 2x^2(x + 4) + 3(x + 4) \\
            & = (x + 4)(2x^2 + 3),
        \end{align*}
        so our equation becomes \((x + 4)(2x^2 + 3) = 0\). Setting each factor equal to 0 gives \(x + 4 = 0\) and \(2x^2 + 3 = 0\). The first equation gives \(x = -4\), but
        the second gives no solution: the quadratic \(2x^2 + 3\) has a minimum value of 3. Alternatively, you could do the algebra:
        \begin{align*}
            2x^2 + 3 & = 0 \\
            2x^2 & = -3 \\
            x^2 & = -\frac{3}{2},
        \end{align*}
        and a square can never be negative.

        \item More factoring! Notice that
        \begin{align*}
            2x^2 + 7x + 3 & = 2x^2 + 6x + x + 3 \\
            & = (2x^2 + 6x) + (x + 3) \\
            & = 2x(x + 3) + 1(x + 3) \\
            & = (x + 3)(2x + 1),
        \end{align*}
        so we can set each factor equal to 0 to solve our original equation. \(x + 3 = 0\) and \(2x + 1 = 0\), so \(x = -3, -\frac{1}{2}\).
    \end{enumerate}

To check your answers when solving any equation, plug them into the original equation, then see if the equation is true.

The last two problems were the more difficult ones. It turns out that in standard form, it's hard to find zeros, but in factored form, the process becomes easier.
\end{explanation}
\end{problem}

\begin{problem}
Ella is working on some math homework and decides to graph the 
function 
\(f\left(x\right)=\left(x-2\right)\left(x+3\right)\left(x-5\right)\left(x+15\right)\) 
on her calculator to get a sense of what the graph looks like. Her 
viewing window on the graphing calculator is set to \([-10,10]\) by \([-10,10]\). 

\begin{enumerate}
    \item How many \(x\)-intercepts will Ella see with this viewing window?
    \item How many \(x\)-intercepts does \(f\) have?
    \item What viewing window would you use?
\end{enumerate}

\begin{explanation}
Remember that \(x\)-intercepts correspond to solutions to the equation \(f(x) = 0\). 
Solving \((x - 2)(x + 3)(x - 5)(x + 15) = 0\), we find \(x = -15, -3, 2, 5\).
\begin{enumerate}
    \item She will see 3 \(x\)-intercepts.
    \item The graph of \(f\) has 4 \(x\)-intercepts.
    \item I would use \([-20, -20]\) by \([-10, 10]\).
\end{enumerate}
\end{explanation}
\end{problem}

\begin{problem}
Consider the following graph of the function \(f\).

\begin{center}
    \includegraphics[width=10cm]{endBehaviorLinear.png}
\end{center}

\begin{enumerate}
    \item Write a sentence describing what happens to \(f\left(x\right)\) 
    as \(x\) gets bigger and bigger in the positive direction. 
    Remember that \(f(x)\) is the \(y\)-value associated with \(x\). 
    An example sentence might be ``As \(x\) gets bigger and bigger in 
    the positive direction, \(f(x)\) gets...''.
    \item Write a sentence describing what happens to \(f\left(x\right)\) 
    as \(x\) gets bigger and bigger in the negative direction. 
\end{enumerate}
\begin{explanation}
    \begin{enumerate}
        \item As \(x\) gets bigger and bigger in the positive direction, \(f(x)\) gets bigger and bigger in the positive direction.
        \item As \(x\) gets bigger and bigger in the negative direction, \(f(x)\) gets bigger and bigger in the negative direction.
    \end{enumerate}
\end{explanation}
\end{problem}

\begin{problem}
Consider the following graph of the function \(f\).

\begin{center}
    \includegraphics[width=10cm]{endBehaviorQuadratic.png}
\end{center}

\begin{enumerate}
    \item Write a sentence describing what happens to \(f\left(x\right)\) 
    as \(x\) gets bigger and bigger in the positive direction. 
    Remember that \(f(x)\) is the \(y\)-value associated with \(x\). 
    An example sentence might be ``As \(x\) gets bigger and bigger in 
    the positive direction, \(f(x)\) gets...''.
    \item Write a sentence describing what happens to \(f\left(x\right)\) 
    as \(x\) gets bigger and bigger in the negative direction. 
\end{enumerate}

\begin{explanation}
    \begin{enumerate}
        \item As \(x\) gets bigger and bigger in the positive direction, \(f(x)\) gets bigger and bigger in the positive direction.
        \item As \(x\) gets bigger and bigger in the negative direction, \(f(x)\) gets bigger and bigger in the positive direction.
    \end{enumerate}
\end{explanation}
\end{problem}

\begin{problem}
Consider the polynomial \(g\left(x\right)=2x^5-5x^4-30x^3+5x^2+88x+60\).

\begin{enumerate}
    \item Identify the degree of the polynomial.
    \item Which of the 6 terms, \(2x^5\), \(5x^4\), \(30x^3\), \(5x^2\),
    \(88x\), or \(60\), is greatest when:
    \begin{enumerate}
        \item \(x = 0\)
        \item \(x = 1\)
        \item \(x = 3\)
        \item \(x = 5\)
    \end{enumerate}
    \item Write a sentence describing what happens to \(g\left(x\right)\) 
    as \(x\) gets bigger and bigger in the positive direction. 
    Remember that \(g(x)\) is the \(y\)-value associated with \(x\). 
    An example sentence might be ``As \(x\) gets bigger and bigger in 
    the positive direction, \(g(x)\) gets...''.
    \item Write a sentence describing what happens to \(g\left(x\right)\) 
    as \(x\) gets bigger and bigger in the negative direction. 
\end{enumerate}
\begin{explanation}
    \begin{enumerate}
        \item The polynomial is degree 5.
        \item This asks you to plug in each \(x\) value to each term and compare them.
        \begin{enumerate}
            \item 60
            \item \(88x\)
            \item \(30x^3\)
            \item \(2x^5\)
        \end{enumerate}
        \item As \(x\) gets bigger and bigger in the positive direction, \(g(x)\) gets bigger and bigger in the positive direction.
        \item As \(x\) gets bigger and bigger in the negative direction, \(g(x)\) gets bigger and bigger in the negative direction.
        
    \end{enumerate}

    Notice that for the last two parts, you can graph the function. 

\end{explanation}
\end{problem}

\end{document}
